% Formatting Guidelines:
% Times New Roman font, 12 point size
% Double spaced, 8-14 pages, including figures and references 
% Initial submission: PDF file

\documentclass[12pt]{article}
\usepackage{times}
\usepackage[margin=1in]{geometry}
\usepackage{setspace}
\usepackage{graphicx}
\usepackage{amsmath}
\usepackage{caption}
\usepackage{booktabs}
\usepackage[hidelinks]{hyperref}

\doublespacing

\begin{document}

\begin{center}
    \textbf{\Large Parameter Invariance in Equity Mean-Reversion Strategies via Out-of-Sample Machine Learning Filtering}
    \vspace{0.5cm}
    
    Jose Miguel Diez Chaupis
\end{center}

\vspace{0.5cm}
\begin{center}
    \textbf{Abstract}
\end{center}

This study investigates the mitigation of look-ahead bias and parameter overfitting in algorithmic trading through the application of a temporally isolated machine learning (ML) filter. Traditional mean-reversion strategies are highly sensitive to manually tuned parameters, such as z-score entry thresholds and signal strength minimums, leading to inflated backtest performance that fails to persist out-of-sample. To address this, an algorithmic trading system was developed integrating classical indicators (z-score, RSI, Bollinger Bands) with a LightGBM classification layer. Crucially, the ML model was trained on an explicitly separated five-year historical window and validated using a strictly out-of-sample two-year period, alongside expanding-window adaptive profiles for risk management. A comprehensive parameter sweep evaluated thirty configurations of entry thresholds and signal strengths. The results demonstrated complete parameter invariance: all thirty configurations yielded identical out-of-sample performance, generating exactly eight high-conviction trades with an average Sharpe ratio of 0.8739 and an average return of 30.49\%. This indicates that an appropriately isolated ML layer effectively supersedes manual parameter optimization by aggressively filtering statistically insignificant signals regardless of the underlying heuristics. These findings suggest that robust ML filtering can insulate trading systems from the fragility of heuristic parameter tuning, presenting a generalized methodology for developing resilient quantitative models.

\newpage
\section*{Introduction}

Algorithmic trading systems frequently employ mean-reversion strategies predicated on the statistical tendency of asset prices to revert to their historical averages following significant deviations. While computationally straightforward, these heuristic models are notoriously susceptible to parameter overfitting. Traders optimize thresholds—such as z-scores for price deviation and Relative Strength Index (RSI) bounds—based on historical data, maximizing backtest performance. However, this process often introduces look-ahead bias and results in systems that perform poorly in out-of-sample, live-market conditions. The inherent fragility of these hardcoded parameters presents a critical challenge in quantitative finance.

Recent advancements in machine learning (ML) offer potential solutions for evaluating trade probability dynamically rather than relying solely on static thresholds. Yet, if ML models are trained on the same data window used for performance backtesting, look-ahead bias is merely shifted from manual tuning to the learning algorithm. 

This study addresses the problem of parameter fragility and look-ahead bias by proposing a temporally isolated, hybrid trading architecture. The system combines classical mean-reversion heuristics with an out-of-sample LightGBM classification layer. The central hypothesis is that an ML filter, when strictly trained on historical data separate from the backtest evaluation period, will induce parameter invariance. By learning the underlying statistical edge of trade setups independently, the model should filter noisy, low-probability signals irrespective of the initial, manually set entry thresholds. This research contributes to the field by demonstrating a methodology for constructing resilient trading systems that are immune to standard forms of in-sample optimization.

\section*{Methods}

\textbf{Data Acquisition and Preparation}

Historical end-of-day price and volume data (Open, High, Low, Close, Volume) were sourced for an expanded universe of twenty equities. These equities were selected via rigorous quantitative screening, targeting mean-reverting characteristics such as a Hurst Exponent below 0.5, a 5-day variance ratio below 1.0, and a 5-year empirical dip-recovery rate exceeding 80\%. Data were fetched programmatically using the `yfinance` library. To prevent look-ahead bias, the data were partitioned temporally into a five-year training block and a subsequent two-year out-of-sample backtesting block. 

\textbf{Feature Engineering and Strategy Generation}

A base mean-reversion strategy was implemented utilizing multiple technical indicators. Primary entry signals were generated when an asset's price deviated negatively beyond a defined z-score threshold. The z-score at time $t$, $z_t$, is calculated as:
\begin{equation}
z_t = \frac{P_t - \mu_{20,t}}{\sigma_{20,t}}
\end{equation}
where $P_t$ is the closing price, $\mu_{20,t}$ is the 20-day simple moving average, and $\sigma_{20,t}$ is the 20-day rolling standard deviation. These signals were subject to a five-layer filtering mechanism, including a trend filter (200-day Simple Moving Average), a volatility regime filter (20-day volatility within the 20th to 80th percentiles), and a distance-from-fair-value filter.

An `IndicatorEngine` computed a comprehensive feature set for each daily observation, including RSI, MACD histogram, Bollinger Band bandwidth, and short-term return lags. Derived features designed to capture signal quality, such as the rate of change of the z-score and the Bollinger Band squeeze percentile, were also calculated. 

\textbf{Machine Learning Filter Implementation}

A LightGBM classifier was employed to predict the probability of a signal resulting in a profitable trade. Labels were assigned based on forward risk-adjusted returns over a five-day horizon. A trade setup at time $t$ was classified with a positive label ($y_t = 1$) only if the forward return exceeded 1.5 times the recent 20-day realized volatility:
\begin{equation}
y_t = 
\begin{cases} 
1 & \text{if } \left(\frac{P_{t+5}}{P_t} - 1\right) > 1.5 \cdot \sigma_{ret, 20, t} \\ 
0 & \text{otherwise} 
\end{cases}
\end{equation}
where $\sigma_{ret, 20, t}$ is the standard deviation of daily returns over the preceding 20 days.

The model was trained exclusively on the five-year training block. To ensure robustness and prevent data leakage, a purged cross-validation methodology was utilized, incorporating a five-day gap between training and validation folds. Hyperparameters were optimized via grid search on the training data. The resulting model was then frozen. During the out-of-sample backtest, any heuristic signal generated was passed through this frozen model. Signals with an ML confidence score below a defined threshold were overridden and rejected.

\textbf{Risk Management and Adaptive Profiles}

Risk management was handled dynamically using Expanding-Window Adaptive Profiles. Position sizing, stop-loss, and take-profit multipliers (based on the Average True Range) were calculated on a rolling, expanding window basis. For any day $T$, parameters were calibrated using data only up to day $T-1$, effectively nullifying look-ahead bias in risk parameterization. 

Furthermore, to adaptively allocate capital based on prevailing market conditions, a Gaussian Mixture Model (GMM) was utilized for unsupervised regime detection. The probability density function of a three-component GMM evaluating daily market returns $x$ is defined as:
\begin{equation}
p(x) = \sum_{k=1}^{3} \pi_k \mathcal{N}(x \mid \mu_k, \Sigma_k)
\end{equation}
where $\pi_k$ represents the mixing weight, and $\mathcal{N}(x \mid \mu_k, \Sigma_k)$ is the multivariate normal distribution for regime $k$ with mean $\mu_k$ and covariance $\Sigma_k$. Capital allocation and position sizing multipliers were dynamically adjusted based on the highest-probability inferred regime.

\textbf{Market-Neutral Pairs Trading}

In addition to directional mean-reversion, the system incorporated a market-neutral pairs trading module. To identify statistically robust pairs, the Engle-Granger two-step method for cointegration was applied. For two asset price series $P^{(A)}_t$ and $P^{(B)}_t$, an Ordinary Least Squares (OLS) regression is first performed:
\begin{equation}
P^{(A)}_t = \alpha + \beta P^{(B)}_t + \varepsilon_t
\end{equation}
where $\beta$ is the hedge ratio. Subsequently, the stationarity of the residuals $\varepsilon_t$ is tested using an Augmented Dickey-Fuller (ADF) test. If the residuals are stationary (i.e., $p$-value $< 0.05$), the pair is deemed cointegrated, allowing the strategy to trade mean-reverting deviations in the spread $\varepsilon_t$.

\textbf{Parameter Sweep and Evaluation}

To test the hypothesis of parameter invariance, a programmatic parameter sweep was conducted over the two-year out-of-sample period. The sweep iterated through thirty unique combinations of two fundamental heuristic parameters: the `Z\_SCORE\_ENTRY\_THRESHOLD` (ranging from 1.3 to 1.8) and the `MIN\_SIGNAL\_STRENGTH` (ranging from 0.0 to 0.30). For each combination, the identical, frozen ML filter and expanding-window adaptive profiles were applied. System performance was evaluated based on the total number of executed trades, total return, Sharpe ratio, and alpha generated against a buy-and-hold SPY benchmark.

\section*{Results}

The parameter sweep tested thirty distinct configurations of the baseline mean-reversion heuristics across the two-year out-of-sample testing window. The configurations ranged from highly permissive (Z-score threshold of 1.3, minimum signal strength of 0.0) to restrictive (Z-score threshold of 1.8, minimum signal strength of 0.30).

The results of the sweep revealed complete uniformity across all tested configurations. Every parameter combination produced identical performance metrics. Specifically, regardless of the initial heuristic parameters, the trading system executed exactly 65 trades over the two-year period. 

Across the expanded 20-stock universe, the strategy yielded an average total return of 38.37\% (17.78\% annualized). The risk-adjusted performance, measured by the Sharpe ratio, was consistent at 0.9915 for all configurations. The strategy generated an average alpha of -0.0688 against the SPY benchmark. 

The LightGBM model, trained on the preceding five-year window, demonstrated a cross-validation score of 0.6241. Feature importance analysis from the training phase indicated that the volatility regime (`vol_regime`), Bollinger Band bandwidth (`bb_bandwidth`), and distance from the 200-day Simple Moving Average (`dist_sma200`) were the most significant predictive features. 

\section*{Discussion}

This study sought to evaluate whether a temporally isolated machine learning filter could mitigate parameter overfitting in a mean-reversion trading strategy. The hypothesis was that the ML model would learn the true statistical edge of the setups and override manual parameter tuning.

The results strongly support this hypothesis. The observation that thirty distinct, manually defined parameter configurations resulted in the exact same 65 trades is highly significant. It demonstrates the phenomenon of parameter invariance induced by the ML layer. When the underlying heuristics were loosened (e.g., setting the z-score threshold to 1.3), the system inherently generated a larger volume of "noisy," lower-probability signals. However, because the ML model was trained on a rigorous five-year out-of-sample dataset, it recognized that these marginal setups lacked a statistical edge and systematically rejected them. Conversely, it consistently identified and approved the same core set of high-conviction trades across all test cases.

This invariance indicates that the system is no longer reliant on the precise, and often fragile, tuning of heuristic parameters. By removing look-ahead bias through strict temporal splitting and expanding-window calibration for risk profiles, the strategy's performance (Sharpe ratio of 0.9915) is arguably a more authentic representation of its out-of-sample efficacy than prior in-sample optimizations that yielded artificially higher returns.

A limitation of this study is the relatively small sample size of executed trades (65) within the two-year test window. While the model successfully avoided bad trades, the low frequency suggests that strict ML filtering on stringent risk-adjusted labels may severely restrict capital deployment opportunities.

These findings have important implications for algorithmic trading system design. They suggest that developers can deploy strategies with moderately permissive baseline parameters, relying on a robustly trained ML filter to serve as the ultimate arbiter of trade execution. Future research should involve forward-testing this architecture via paper trading in live market conditions to validate the persistence of this edge and explore techniques to increase trade frequency without compromising the established filtering stringency. 

\section*{Acknowledgements}

This research was conducted using resources provided by the University of California, Berkeley. We thank the developers of the `pandas`, `scikit-learn`, and `LightGBM` libraries.

\section*{References}

[1] Bailey, D. H., Borwein, J. M., Lopez de Prado, M., \& Zhu, Q. J. (2014). Pseudo-mathematics and financial charlatanism: The effects of backtest overfitting on out-of-sample performance. \textit{Notices of the AMS}, 61(5), 458-471.

[2] Lopez de Prado, M. (2018). \textit{Advances in financial machine learning}. John Wiley \& Sons.

[3] Aronson, D. R. (2007). \textit{Evidence-based technical analysis: Applying the scientific method and statistical inference to trading signals}. John Wiley \& Sons.

[4] Ke, G., Meng, Q., Finley, T., Wang, T., Chen, W., Ma, W., ... \& Liu, T. Y. (2017). Lightgbm: A highly efficient gradient boosting decision tree. \textit{Advances in neural information processing systems}, 30.

[5] Wilder, J. W. (1978). \textit{New concepts in technical trading systems}. Trend Research.

[6] Bollinger, J. (2001). \textit{Bollinger on Bollinger bands}. McGraw Hill Professional.

[7] Engle, R. F., \& Granger, C. W. (1987). Co-integration and error correction: representation, estimation, and testing. \textit{Econometrica: journal of the Econometric Society}, 251-276.

[8] Krauss, C., Do, X. A., \& Huck, N. (2017). Deep neural networks, gradient-boosted trees, random forests: Statistical arbitrage on the S\&P 500. \textit{European Journal of Operational Research}, 259(2), 689-702.

[9] Guida, T. (2018). \textit{Big data and machine learning in quantitative investment}. John Wiley \& Sons.

[10] Dixon, M. F., Halperin, I., \& Bilokon, P. (2020). \textit{Machine learning in finance} (Vol. 1170). Berlin: Springer.

[11] Jansen, S. (2020). \textit{Machine learning for algorithmic trading}. Packt Publishing Ltd.

[12] Pardo, R. (2011). \textit{The evaluation and optimization of trading strategies}. John Wiley \& Sons.

[13] Chan, E. (2008). \textit{Quantitative trading: how to build your own algorithmic trading business}. John Wiley \& Sons.

[14] Hull, J. C. (2003). \textit{Options futures and other derivatives}. Pearson Education India.

[15] Pedregosa, F., Varoquaux, G., Gramfort, A., Michel, V., Thirion, B., Grisel, O., ... \& Duchesnay, E. (2011). Scikit-learn: Machine learning in Python. \textit{the Journal of machine Learning research}, 12, 2825-2830.

\end{document}
